\chapter{The CATIA mass/CG skill}
\label{ch:catia}

The skill lets an engineer ask, in the chat UI, for the mass and centre of
gravity of a CATIA assembly, and --- after explicitly approving the preview ---
export the result as an Adams/Car \file{.cmd} file.

It is the one feature in the project that reaches outside the application: it
drives CATIA over COM and writes files to disk. Everything about its design
follows from that.

\section{Where the gates live}

The security gates --- the command allow-list, the placeholder check, the human
approval gate and the channel separation --- are \textbf{not reimplemented} in
the web application. The skill package \file{skill/catia-mass-cg.skill} carries
its own harness (\file{runner/agent.py}), and \code{CatiaSkillService} unpacks
the archive at runtime and imports \code{validate}, \code{dispatch},
\code{Gate} and \code{for\_model} from it.

The consequence is the point: updating the \file{.skill} file updates the gates
with it, from one source. The service adds only what a web application has to
add:

\begin{itemize}
\item unpacking the archive under \code{<DATA\_DIR>/skills}, stamped with the
archive hash so a new build lands in a new folder;
\item session management --- message history, gate state, one workspace per
user;
\item the tool-calling chat request to Ollama, made with \code{httpx} so a slow
model cannot take the server down;
\item carrying screen-channel events (the preview table, the command lines) to
the UI; and
\item a deterministic export path for the approval button, which does not depend
on the model remembering to call \code{export}.
\end{itemize}

\section{Sessions and workspaces}

A session holds the message history and the approval gate, and is kept
\textbf{in memory}. A server restart clears the conversation, but not the work:
runs and \file{memory.sqlite} live in the workspace and come back through
\code{history}.

The workspace is \code{<CATIA\_SKILL\_WORKSPACE\_ROOT>/<username>}, where the
username comes from the auth session --- or is \code{local} when authentication
is off. On first use the workspace is seeded with the example asset files from
the skill package (sub-assembly map, Adams map, transform profile).

\begin{callout}{Warning}
With \env{APP\_AUTH\_ENABLED=false} \textbf{every caller shares the \code{local}
workspace and its run history.} Turn authentication on if several engineers
should keep separate histories.
\end{callout}

\section{Running over the network}

The endpoints accept any client by default, so an engineer can drive the skill
from another machine. Serve on the LAN the usual way:

\begin{shellblock}
$env:UCAK_HOST = "0.0.0.0"
.\start.bat
\end{shellblock}

Two things do not travel with the client, and the UI says so to a remote user:

\begin{itemize}
\item \textbf{The measurement runs on the server's CATIA}, over COM, and the
Adams/Car \file{.cmd} file is written to the \textbf{server's} disk --- the
workspace folder shown in the module header. Enable the skill on the machine
CATIA actually runs on.
\item \textbf{CATIA answers one measurement at a time.} Concurrent \code{cmc}
commands are serialized behind a lock; a second engineer gets ``another
measurement is running'' rather than a corrupted result.
\end{itemize}

To restrict which clients may use the skill, list them --- \code{;} or \code{,}
separated, where \code{local}/\code{localhost} covers every loopback spelling.
Empty means no restriction:

\begin{shellblock}
$env:CATIA_SKILL_ALLOWED_CLIENTS = "local;10.0.0.7"
\end{shellblock}

\section{Endpoints}

\subsection*{\route{GET /skills/catia-mass-cg/status}}

What the UI polls to decide whether to show the module at all.

\begin{jsonblock}
{
  "available": true,
  "enabled": true,
  "source": "fake",
  "model": "qwen3:4b-instruct",
  "ollama_host": "http://127.0.0.1:11434",
  "skill_root": "...\\data\\skills\\catia-mass-cg-<hash>",
  "workspace_root": "...\\data\\catia_skill",
  "sessions": 1,
  "client_allowed": true,
  "remote_client": false
}
\end{jsonblock}

\code{client\_allowed} lets the interface say \emph{why} the module is unusable
instead of presenting an empty chat that will answer \code{403}.
\code{remote\_client} is information, not a block: it is how a remote engineer is
told the measurement will run on the server. With the skill disabled the
response is simply \code{\{"enabled": false\}}.

\subsection*{\route{POST /skills/catia-mass-cg/chat}}

One conversation turn. Send \code{message} and, after the first turn, the
\code{session\_id} from the previous response. Returns the events produced by the
turn.

\subsection*{\route{POST /skills/catia-mass-cg/approve}}

Approves the pending preview and runs the export deterministically through the
harness. The gate is still checked inside the harness --- run-id matching, the
approval code read from disk, \code{E\_STALE\_APPROVAL} raised by \code{cmc} ---
so pressing the button cannot approve something other than what was shown.

\subsection{Events}

A response carries a list of events, each with a \code{kind}.

\begin{longtable}{L{0.16\textwidth} L{0.74\textwidth}}
\toprule
\textbf{Kind} & \textbf{Meaning} \\
\midrule
\endfirsthead
\bottomrule
\endfoot
\code{model} & Text the model produced for the user. \\
\code{command} & A \code{cmc} command the harness accepted and ran. \\
\code{result} & The structured result of that command, with \code{status},
\code{code} and a Turkish \code{message\_tr} / \code{hint\_tr}. \\
\code{screen} & Something to render: the preview table, command output. \\
\code{harness} & A note from the harness itself, such as a nudge after a
malformed tool call. \\
\end{longtable}

Alongside the events, the response carries \code{state},
\code{approval\_pending} and \code{pending\_run\_id} --- which is what the UI
needs to decide whether to show the approval button.

\subsection{Status codes}

\begin{longtable}{L{0.12\textwidth} L{0.78\textwidth}}
\toprule
\textbf{Code} & \textbf{When} \\
\midrule
\endfirsthead
\bottomrule
\endfoot
\code{404} & The skill is disabled (\env{CATIA\_SKILL\_ENABLED=false}), or the
session id is unknown or expired. \\
\code{403} & The client is not on \env{CATIA\_SKILL\_ALLOWED\_CLIENTS}. \\
\code{409} & A request is already running in this session, or there is nothing
pending to approve. \\
\code{503} & The skill package or the model is unavailable. \\
\end{longtable}

\section{Practice mode}

\env{CATIA\_SKILL\_SOURCE=fake} --- the default --- drives a synthetic assembly
shipped with the skill. No CATIA, no COM, no licence: the whole flow, including
the approval gate and the export, can be rehearsed on any machine. Set it to
\code{catia} deliberately, on the workstation with CATIA, when the measurements
should be real.

\section{Operational checks}

\begin{shellblock}
# from the server itself
Invoke-RestMethod http://127.0.0.1:8000/skills/catia-mass-cg/status

# from another machine on the LAN -- expect remote_client: true
Invoke-RestMethod http://SERVER-IP:8000/skills/catia-mass-cg/status
\end{shellblock}

The model needs reliable tool calling: it has to emit well-formed tool calls turn
after turn, because a malformed one costs a nudge and the harness only allows
\env{CATIA\_SKILL\_MAX\_NUDGES} of them. \code{qwen3:4b-instruct} at a
Q5\_K\_M quantization or better is the tested configuration. A smaller or more
aggressively quantized model tends to narrate the command instead of calling it,
which the harness rejects rather than guesses at.
